% ============================================================
\section{Ecole Reward Mechanism and the One-Step Solve}

This section explains in depth how the \texttt{OptimalSol} reward class and \texttt{solve\_instance} function interact with Ecole internals, and why a single call to \texttt{env.step()} is sufficient to retrieve the optimal objective value of a MIP instance.

% ============================================================
\subsection{The Ecole Environment Interface}

Ecole models combinatorial optimization as a sequential decision process following the OpenAI Gym convention. Every Ecole environment exposes two primary methods:

\begin{itemize}
    \item \texttt{env.reset(instance\_path)} --- loads a new MIP instance into SCIP, initializes internal state, and returns the first observation.
    \item \texttt{env.step(action)} --- applies one \emph{action} (e.g., a branching decision or a parameter configuration), advances the solver, and returns a 5-tuple \texttt{(observation, action\_set, reward, done, info)}.
\end{itemize}

The \texttt{reward} in the returned tuple is not computed inside \texttt{env.step()} itself. Instead, Ecole delegates reward computation to a user-supplied \emph{reward function} object --- in this script, an instance of \texttt{OptimalSol}.

% ============================================================
\subsection{What Is \texttt{extract} and When Is It Called?}

Every Ecole reward function must implement two methods that Ecole calls automatically at specific points in the episode lifecycle:

\begin{lstlisting}
class OptimalSol:
    def before_reset(self, model): pass

    def extract(self, model, done=False):
        pyscipopt_model = model.as_pyscipopt()
        if done:
            reward = pyscipopt_model.getObjVal(original=True)
        else:
            reward = None
        return reward
\end{lstlisting}

\begin{description}
    \item[\texttt{before\_reset(model)}]
        Called by Ecole \emph{before} \texttt{env.reset()} loads a new instance. It is the place to initialize any per-episode state. Here it is a no-op because \texttt{OptimalSol} holds no state of its own.

    \item[\texttt{extract(model, done)}]
        Called by Ecole \emph{inside every} \texttt{env.step()} call, \emph{after} the solver has advanced by one action. Ecole passes two arguments:
        \begin{itemize}
            \item \texttt{model} --- the live Ecole model wrapper around the current SCIP state.
            \item \texttt{done} --- a boolean that is \texttt{True} only when the episode has ended (i.e.\ SCIP has finished solving and there are no more decisions to make).
        \end{itemize}
        Whatever value \texttt{extract} returns becomes the \texttt{reward} field in \texttt{env.step()}'s return tuple.
\end{description}

The conditional logic inside \texttt{extract} is therefore:
\begin{itemize}
    \item If \texttt{done=False} (solver still running, mid-episode): return \texttt{None} --- no meaningful reward yet.
    \item If \texttt{done=True} (solver finished): call \texttt{pyscipopt\_model.getObjVal(original=True)} and return the optimal objective value.
\end{itemize}

The flag \texttt{original=True} in \texttt{getObjVal} instructs PySCIPOpt to report the value in the \emph{original} objective space, undoing any internal transformations SCIP may have applied (e.g., negating a maximisation objective to convert it to minimisation form).

% ============================================================
\subsection{Why One \texttt{env.step()} Is Enough}

The key is the choice of Ecole environment: \texttt{ecole.environment.Configuring}.

\subsubsection{Configuring vs.\ Branching Environments}

Most Ecole environments (e.g., \texttt{Branching}) model branch-and-bound as a \emph{multi-step} process: each call to \texttt{env.step()} applies a single branching decision and returns control to Python, so the episode runs for as many steps as there are branching nodes. These environments are used for training branching policies.

\texttt{Configuring} is different. Its action space is a dictionary of SCIP parameters to set before the solve begins. Once the parameters are applied, SCIP runs \emph{to completion} --- the entire branch-and-bound tree is explored internally, without returning to Python at intermediate nodes. The episode therefore consists of exactly \textbf{one step}:

\begin{enumerate}
    \item \texttt{env.reset(instance\_path)} --- loads the instance; \texttt{before\_reset} is called.
    \item \texttt{env.step(\{\})} --- applies no parameter changes (empty dict), triggers SCIP to solve the full MIP, and returns once the solve is complete with \texttt{done=True}.
\end{enumerate}

\subsubsection{The Internal Call Sequence}

When \texttt{env.step(\{\})} is called in \texttt{solve\_instance}:

\begin{enumerate}
    \item Ecole applies the (empty) parameter dictionary to SCIP --- a no-op here.
    \item SCIP runs its full branch-and-bound algorithm internally until it finds and verifies the optimal solution.
    \item Ecole sets \texttt{done = True} and calls \texttt{reward\_fun.extract(model, done=True)}.
    \item \texttt{extract} queries \texttt{pyscipopt\_model.getObjVal(original=True)} and returns the optimal value.
    \item \texttt{env.step()} returns \texttt{(None, None, optimal\_value, True, \{\})}.
\end{enumerate}

Back in \texttt{solve\_instance}, only the third element of the tuple is used:

\begin{lstlisting}
_, _, solution, _, _ = env.step({})
out_queue.put({instance: solution})
\end{lstlisting}

The variables \texttt{observation}, \texttt{action\_set}, \texttt{done}, and \texttt{info} are all discarded with \texttt{\_} because they are irrelevant for pure instance solving.

\subsubsection{Summary of the Episode Lifecycle}

\begin{center}
\begin{tabular}{lll}
\toprule
\textbf{Call} & \textbf{Ecole internally calls} & \textbf{Result} \\
\midrule
\texttt{env.reset(path)} & \texttt{before\_reset(model)} & instance loaded into SCIP \\
\texttt{env.step(\{\})} & SCIP solves full MIP & \texttt{done=True} \\
                        & \texttt{extract(model, done=True)} & returns \texttt{ObjVal} \\
\bottomrule
\end{tabular}
\end{center}

Because the \texttt{Configuring} environment always terminates in one step, \texttt{extract} is only ever called once per instance with \texttt{done=True}, so the \texttt{done=False} branch (which returns \texttt{None}) is never actually reached in this script. It exists only to satisfy the Ecole reward function interface, which is designed to also support multi-step environments.

\subsection{Contrast With Training Usage}

In the reinforcement learning training phase (e.g., \texttt{05\_train\_rl.py}), a \emph{branching} environment is used instead. There, \texttt{env.step()} is called once per branching decision, and \texttt{extract} is called with \texttt{done=False} at every intermediate node and with \texttt{done=True} only at the end. A reward function in that context might return incremental values at each step. The \texttt{OptimalSol} class here is a degenerate case of that general interface, repurposed for simple instance solving with no learning.
